\subsection{Variation 1.1}
\textbf{Question:} Which of the following are roles of an Operating System (OS)? Please select the \textbf{two} correct answers.\\
\begin{tabularx}{\textwidth}{|X|X|}
    \hline
    \textbf{Correct Answer(s)} & \textbf{Incorrect Answers} \\
    \hline
    Managing hardware resources. & Designing hardware components like CPUs and GPUs. \\
    Acting as a resource manager for CPU, memory, and storage. & Enhancing the speed of computer hardware. \\
        & Ensuring the physical security of the computer system. \\
    \hline
\end{tabularx}
\textbf{Solution \& Feedback:}
\begin{itemize}
    \item \textbf{Managing hardware resources:} The OS manages hardware resources such as the CPU, memory, and I/O devices to ensure they are allocated effectively to running programs.
    \item \textbf{Acting as a resource manager for CPU, memory, and storage:} One of the primary roles of an OS is to act as a resource manager, ensuring efficient use of system resources.
    \item \textbf{Designing hardware components like CPUs and GPUs:} The OS is software and does not design hardware. Hardware design is the responsibility of engineers and manufacturers.
    \item \textbf{Enhancing the speed of computer hardware:} The OS can optimize resource utilization but cannot physically increase hardware speed, which is determined by the hardware's design.
    \item \textbf{Ensuring the physical security of the computer system:} The OS can enforce software-level security, such as user authentication and encryption, but physical security is handled by external measures like locks or surveillance.
\end{itemize}

\subsection{Variation 1.2}
\begin{itemize}
    \item \textbf{Question:} Which of the following are roles of an Operating System (OS)? Please select the two correct answers.\\
    \begin{tabularx}{\textwidth}{|X|X|}
        \hline
        \textbf{Correct Answer(s)} & \textbf{Incorrect Answers} \\
        \hline
        Enabling software applications to run efficiently. & Enhancing the speed of computer hardware. \\
        Providing an interface between users and hardware & Preventing hardware from overheating. \\
         & Ensuring the physical security of the computer system. \\
        \hline
    \end{tabularx}
    
    \item \textbf{Solution \& Feedback:}
    \begin{itemize}
        \item \textbf{Enabling software applications to run efficiently:} The OS provides necessary services, such as memory management and scheduling, to ensure software applications run smoothly.
        \item \textbf{Providing an interface between users and hardware:} The OS provides an abstraction layer that simplifies user interaction with the underlying hardware.
        \item \textbf{Enhancing the speed of computer hardware:} The OS can optimize resource utilization but cannot physically increase hardware speed, which is determined by the hardware's design.
        \item \textbf{Preventing hardware from overheating:} While the OS may monitor hardware temperatures and trigger system shutdowns to prevent damage, actual prevention of overheating depends on physical components like cooling systems.
        \item \textbf{Ensuring the physical security of the computer system:} The OS can enforce software-level security, such as user authentication and encryption, but physical security is handled by external measures like locks or surveillance.
    \end{itemize}
\end{itemize}

\subsection{Variation 1.3}
\begin{itemize}
    \item \textbf{Question:} Which of the following are roles of an Operating System (OS)? Please select the \textbf{two} correct answers.
    
    \begin{tabularx}{\textwidth}{|X|X|}
        \hline
        \textbf{Correct Answer(s)} & \textbf{Incorrect Answers} \\
        \hline
        Managing hardware resources. & Enhancing the speed of computer hardware. \\
        Enabling software applications to run efficiently. & Ensuring the physical security of the computer system. \\
         & Designing hardware components like CPUs and GPUs. \\
        \hline
    \end{tabularx}
    
    \item \textbf{Solution \& Feedback:}
    \begin{itemize}
        \item \textbf{Managing hardware resources:} The OS manages hardware resources such as the CPU, memory, and I/O devices to ensure they are allocated effectively to running programs.
        \item \textbf{Enabling software applications to run efficiently:} The OS provides necessary services, such as memory management and scheduling, to ensure software applications run smoothly.
        \item \textbf{Enhancing the speed of computer hardware:} The OS can optimize resource utilization but cannot physically increase hardware speed, which is determined by the hardware's design.
        \item \textbf{Ensuring the physical security of the computer system:} The OS can enforce software-level security, such as user authentication and encryption, but physical security is handled by external measures like locks or surveillance.
        \item \textbf{Designing hardware components like CPUs and GPUs:} The OS is software and does not design hardware. Hardware design is the responsibility of engineers and manufacturers.
    \end{itemize}
\end{itemize}

\subsection{Variation 2.1}
\begin{itemize}
    \item \textbf{Question:} Which of the following statements are true about algorithm complexity? Please select the two correct answers.
    
    \begin{tabularx}{\textwidth}{|X|X|}
        \hline
        \textbf{Correct Answer(s)} & \textbf{Incorrect Answers} \\
        \hline
        Algorithm complexity measures the efficiency of an algorithm in terms of time and space. & Big-O notation is used to describe the exact number of operations in an algorithm. \\
        Time complexity often focuses on the worst-case scenario for large input sizes. & The complexity class O(1) represents logarithmic growth. \\
         & Algorithm complexity is only relevant for programs implemented in a specific programming language. \\
        \hline
    \end{tabularx}
    
    \item \textbf{Solution \& Feedback:}
    \begin{itemize}
        \item \textbf{Algorithm complexity measures the efficiency of an algorithm in terms of time and space:} Algorithm complexity assesses both the time required and the memory used by an algorithm as functions of input size.
        \item \textbf{Time complexity often focuses on the worst-case scenario for large input sizes:} Worst-case analysis ensures that the algorithm performs acceptably even in the most demanding scenarios, which is critical for reliability.
        \item \textbf{Big-O notation is used to describe the exact number of operations in an algorithm:} Big-O notation describes the growth rate or complexity class, not the exact number of operations.
        \item \textbf{The complexity class O(1) represents logarithmic growth:} O(1) represents constant time complexity, meaning the algorithm's performance does not depend on the input size.
        \item \textbf{Algorithm complexity is only relevant for programs implemented in a specific programming language:} Algorithm complexity is independent of programming languages and measures the theoretical efficiency of the algorithm itself, not its implementation.
    \end{itemize}
\end{itemize}

\subsection{Variation 2.2}
\begin{itemize}
    \item \textbf{Question:} Which of the following statements are true about algorithm complexity? Please select the two correct answers.
    
    \begin{tabularx}{\textwidth}{|X|X|}
        \hline
        \textbf{Correct Answer(s)} & \textbf{Incorrect Answers} \\
        \hline
        O(n\textsuperscript{2}) describes an algorithm where the number of operations grows quadratically with input size. & Big-O notation is used to describe the exact number of operations in an algorithm. \\
        In algorithm design, the programmer often needs to balance considerations of time and space complexity. & The complexity class O(1) represents logarithmic growth. \\
         & Algorithm complexity is only relevant for programs implemented in a specific programming language. \\
        \hline
    \end{tabularx}
    
    \item \textbf{Solution \& Feedback:}
    \begin{itemize}
        \item \textbf{O(n\textsuperscript{2}) describes an algorithm where the number of operations grows quadratically with input size:} This is the definition of quadratic complexity, common in algorithms like bubble sort.
        \item \textbf{In algorithm design, the programmer often needs to balance considerations of time and space complexity:} For example, using extra memory (space) for caching can reduce the number of operations (time), and vice versa.
        \item \textbf{Big-O notation is used to describe the exact number of operations in an algorithm:} Big-O notation describes the growth rate or complexity class, not the exact number of operations.
        \item \textbf{The complexity class O(1) represents logarithmic growth:} O(1) represents constant time complexity, meaning the algorithm's performance does not depend on the input size.
        \item \textbf{Algorithm complexity is only relevant for programs implemented in a specific programming language:} Algorithm complexity is independent of programming languages and measures the theoretical efficiency of the algorithm itself, not its implementation.
    \end{itemize}
\end{itemize}

\subsubsection{Variation 2.3}
\begin{itemize}
    \item \textbf{Question:} Which of the following statements are true about algorithm complexity? Please select the \textbf{two} correct answers.
    
    \begin{tabularx}{\textwidth}{|X|X|}
        \hline
        \textbf{Correct Answer(s)} & \textbf{Incorrect Answers} \\
        \hline
        Exponential complexity (O(2\textsuperscript{n})) is generally considered intractable for large inputs. & Big-O notation is used to describe the exact number of operations in an algorithm. \\
        Time complexity often focuses on the worst-case scenario for large input sizes. & The complexity class O(1) represents logarithmic growth. \\
         & Algorithm complexity is only relevant for programs implemented in a specific programming language. \\
        \hline
    \end{tabularx}
    
    \item \textbf{Solution \& Feedback:}
    \begin{itemize}
        \item \textbf{Exponential complexity (O(2\textsuperscript{n})) is generally considered intractable for large inputs:} Exponential algorithms grow so rapidly with input size that they become impractical for anything beyond small inputs.
        \item \textbf{Time complexity often focuses on the worst-case scenario for large input sizes:} Worst-case analysis ensures that the algorithm performs acceptably even in the most demanding scenarios, which is critical for reliability.
        \item \textbf{Big-O notation is used to describe the exact number of operations in an algorithm:} Big-O notation describes the growth rate or complexity class, not the exact number of operations.
        \item \textbf{The complexity class O(1) represents logarithmic growth:} O(1) represents constant time complexity, meaning the algorithm's performance does not depend on the input size.
        \item \textbf{Algorithm complexity is only relevant for programs implemented in a specific programming language:} Algorithm complexity is independent of programming languages and measures the theoretical efficiency of the algorithm itself, not its implementation.
    \end{itemize}
\end{itemize}

\subsubsection{Variation 3.1}
\begin{itemize}
    \item \textbf{Question:} Which of the following statements about interrupts are true? Please select the three correct answers.
    
    \begin{tabularx}{\textwidth}{|X|X|}
        \hline
        \textbf{Correct Answer(s)} & \textbf{Incorrect Answers} \\
        \hline
        Interrupts can be triggered by both hardware and software & Interrupts cannot be triggered by software \\
        The hardware triggers interrupts by sending a signal to the CPU & Hardware can only trigger interrupts at specific times \\
        Hardware can trigger an interrupt at any time & \\
        \hline
    \end{tabularx}
   \\ 
    \item \textbf{Solution \& Feedback:}
    \begin{itemize}
        \item Interrupts can be triggered by hardware (for example, a mouse click) or software (for example, when a program tries to access a file, which leads to a system call).
        \item Hardware may trigger an interrupt at any time by sending a signal to the CPU.
    \end{itemize}
\end{itemize}

\subsection{Variation 3.2}
\begin{itemize}
    \item \textbf{Question:} Which of the following statements about interrupts are true? Please select the three correct answers.
    
    \begin{tabularx}{\textwidth}{|X|X|}
        \hline
        \textbf{Correct Answer(s)} & \textbf{Incorrect Answers} \\
        \hline
        Interrupts can be triggered by both hardware and software & Interrupts cannot be triggered by hardware \\
        Software triggers interrupts by executing a special instruction & Interrupts are triggered in kernel mode only \\
        System calls are invoked using software interrupts & \\
        \hline
    \end{tabularx}
    
    \item \textbf{Solution \& Feedback:}
    \begin{itemize}
        \item Interrupts can be triggered by hardware (for example, a mouse click) or software (for example, when a program tries to access a file, which leads to a system call).
        \item Software triggers an interrupt by executing a special operation – to make a system call and switch to kernel mode.
    \end{itemize}
\end{itemize}

\subsection{Variation 3.3}
\begin{itemize}
    \item \textbf{Question:} Which of the following statements about interrupts are true? Please select the three correct answers.
    
    \begin{tabularx}{\textwidth}{|X|X|}
        \hline
        \textbf{Correct Answer(s)} & \textbf{Incorrect Answers} \\
        \hline
        Interrupts can be triggered by both hardware and software & When an interrupt occurs, the CPU waits until the current process is completed before transferring control to the interrupt service routine \\
        When an interrupt occurs, the CPU immediately stops what it is doing and transfers control to the interrupt service routine & Hardware can only trigger interrupts at specific times \\
        Hardware devices can trigger an interrupt at any time & \\
        \hline
    \end{tabularx}
    
    \item \textbf{Solution \& Feedback:}
    \begin{itemize}
        \item The CPU reacts to an interrupt by stopping what it is doing and immediately transferring execution to a fixed location (usually containing the starting address of the interrupt service routine).
        \item Hardware may trigger an interrupt at any time by sending a signal to the CPU.
    \end{itemize}
\end{itemize}

\subsection{Variation 4.1}
\begin{itemize}
    \item \textbf{Question:} For which ones of the following tasks, we can define an algorithm? Please select the three correct answers.
    
    \begin{tabular}{|l|l|}
        \hline
        \textbf{Correct Answer(s)} & \textbf{Incorrect Answers} \\
        \hline
        Searching for an Element & Predicting the Future Human Behaviour \\
        Finding the Greatest Common Divisor & Defining Artistic Creativity \\
        Navigating a Robot in a Grid with Obstacles & \\
        \hline
    \end{tabular}
    
    \item \textbf{Solution \& Feedback:}
    \begin{itemize}
        \item \textbf{Searching for an Element:} Find the position of a specific value in a data structure. Example: Binary Search, Linear Search.
        \item \textbf{Finding the Greatest Common Divisor:} Algorithm: Euclidean Algorithm.
        \item \textbf{Navigating a Robot in a Grid with Obstacles:} Algorithm: A* Pathfinding.
        \item \textbf{Predicting the Future Human Behaviour:} Human actions are influenced by emotions, free will, and other unpredictable factors, making them non-algorithmic.
        \item \textbf{Defining Artistic Creativity:} Creativity is subjective and context-dependent, making it hard to encapsulate in a deterministic algorithm.
    \end{itemize}
\end{itemize}

\subsection{Variation 4.2}
\begin{itemize}
    \item \textbf{Question:} For which ones of the following tasks, we can define an algorithm? Please select the three correct answers.
    
    \begin{tabular}{|l|l|}
        \hline
        \textbf{Correct Answer(s)} & \textbf{Incorrect Answers} \\
        \hline
        Finding the Shortest Path & Understanding Human Consciousness \\
        Solving Systems of Linear Equations & Determining the Morality of Actions \\
        Recognizing Handwritten Digits & \\
        \hline
    \end{tabular}
    
    \item \textbf{Solution \& Feedback:}
    \begin{itemize}
        \item \textbf{Finding the Shortest Path:} Calculate the shortest path between two nodes in a graph. Example: Dijkstra's Algorithm, A* Algorithm.
        \item \textbf{Solving Systems of Linear Equations:} Algorithms: Gaussian Elimination, LU Decomposition.
        \item \textbf{Recognizing Handwritten Digits:} Algorithm: Convolutional Neural Networks (CNNs).
        \item \textbf{Understanding Human Consciousness:} There is no algorithm to fully model or understand subjective experiences and consciousness.
        \item \textbf{Determining the Morality of Actions:} Morality involves subjective cultural, social, and personal values, which cannot be universally defined or computed.
    \end{itemize}
\end{itemize}

\subsection{Variation 4.3}
\begin{itemize}
    \item \textbf{Question:} For which ones of the following tasks, we can define an algorithm? Please select the three correct answers.
    
    \begin{tabular}{|l|l|}
        \hline
        \textbf{Correct Answer(s)} & \textbf{Incorrect Answers} \\
        \hline
        Solving a Maze Puzzle & Defining Universal Humor \\
        Calculating Factorials for Large Numbers & Fully Translating Poetry Between Languages \\
        Face Detection in Images & \\
        \hline
    \end{tabular}
    
    \item \textbf{Solution \& Feedback:}
    \begin{itemize}
        \item \textbf{Solving a Maze Puzzle:} Navigate through a maze to find the exit. Example: Depth-First Search (DFS), Breadth-First Search (BFS).
        \item \textbf{Calculating Factorials for Large Numbers:} Algorithm: Iterative or recursive multiplication, or using modular arithmetic for efficiency.
        \item \textbf{Face Detection in Images:} Algorithms: Viola-Jones Framework, Haar Cascades.
        \item \textbf{Defining Universal Humor:} Humor is context-sensitive and depends on cultural, emotional, and individual preferences, which cannot be universally quantified.
        \item \textbf{Fully Translating Poetry Between Languages:} Poetry involves cultural nuances, emotions, and artistic interpretation that cannot be perfectly captured by deterministic rules.
    \end{itemize}
\end{itemize}

\subsection{Variation 5.1}
\begin{itemize}
    \item \textbf{Question:} Consider the following processes with their burst times and arrival times:
    
    \begin{tabular}{|l|l|l|}
        \hline
        \textbf{Process} & \textbf{Arrival Time (ms)} & \textbf{Burst Time (ms)} \\
        \hline
        P1 & 0 & 5 \\
        P2 & 2 & 3 \\
        P3 & 4 & 8 \\
        P4 & 26 & 6 \\
        \hline
    \end{tabular}
    
    Using the \textbf{First-Come-First-Served (FCFS)} scheduling algorithm, what is the \textbf{average waiting time}?
    
    \begin{tabular}{|l|l|}
        \hline
        \textbf{Correct Answer(s)} & \textbf{Incorrect Answers} \\
        \hline
        1.75 & 5.5 \\
         & 2.5 \\
         & 3 \\
         & 4.25 \\
        \hline
    \end{tabular}
    
    \item \textbf{Solution \& Feedback:}
    \begin{itemize}
        \item \textbf{Step 1: Calculate the finish time for each process:}
        \begin{align*}
            \text{Finish Time (P1)} &= \text{Arrival Time (P1)} + \text{Burst Time (P1)} = 0 + 5 = 5 \\
            \text{Finish Time (P2)} &= \max(\text{Finish Time (P1)}, \text{Arrival Time (P2)}) + \text{Burst Time (P2)} = 5 + 3 = 8 \\
            \text{Finish Time (P3)} &= \max(\text{Finish Time (P2)}, \text{Arrival Time (P3)}) + \text{Burst Time (P3)} = 8 + 8 = 16 \\
            \text{Finish Time (P4)} &= \max(\text{Finish Time (P3)}, \text{Arrival Time (P4)}) + \text{Burst Time (P4)} = 26 + 6 = 32 \\
        \end{align*}
        
        \item \textbf{Step 2: Calculate the waiting time (W) for each process: W = Start Time – Arrival Time}
        \begin{align*}
            W1 &= 0 - 0 = 0 \\
            W2 &= 5 - 2 = 3 \\
            W3 &= 8 - 4 = 4 \\
            W4 &= 26 - 26 = 0 \\
        \end{align*}
        
        \item \textbf{Step 3: Calculate the average:}
        \[
            \text{Average Waiting Time} = \frac{0 + 3 + 4 + 0}{4} = 1.75
        \]
    \end{itemize}
\end{itemize}

\subsection{Variation 5.2}
\begin{itemize}
    \item \textbf{Question:} Consider the following processes with their burst times and arrival times:
    
    \begin{tabular}{|l|l|l|}
        \hline
        \textbf{Process} & \textbf{Arrival Time (ms)} & \textbf{Burst Time (ms)} \\
        \hline
        P1 & 0 & 3 \\
        P2 & 1 & 5 \\
        P3 & 5 & 6 \\
        P4 & 16 & 9 \\
        \hline
    \end{tabular}
    
    Using the \textbf{First-Come-First-Served (FCFS)} scheduling algorithm, what is the \textbf{average waiting time}?
    
    \begin{tabular}{|l|l|}
        \hline
        \textbf{Correct Answer(s)} & \textbf{Incorrect Answers} \\
        \hline
        1.25 & 5.75 \\
         & 3 \\
         & 2.75 \\
         & 3.25 \\
        \hline
    \end{tabular}
    
    \item \textbf{Solution \& Feedback:}
    \begin{itemize}
        \item \textbf{Step 1: Calculate the finish time for each process:}
        \begin{align*}
            \text{Finish Time (P1)} &= \text{Arrival Time (P1)} + \text{Burst Time (P1)} = 0 + 3 = 3 \\
            \text{Finish Time (P2)} &= \max(\text{Finish Time (P1)}, \text{Arrival Time (P2)}) + \text{Burst Time (P2)} = 3 + 5 = 8 \\
            \text{Finish Time (P3)} &= \max(\text{Finish Time (P2)}, \text{Arrival Time (P3)}) + \text{Burst Time (P3)} = 8 + 6 = 14 \\
            \text{Finish Time (P4)} &= \max(\text{Finish Time (P3)}, \text{Arrival Time (P4)}) + \text{Burst Time (P4)} = 16 + 9 = 25 \\
        \end{align*}
        
        \item \textbf{Step 2: Calculate the waiting time (W) for each process: W = Start Time – Arrival Time}
        \begin{align*}
            W1 &= 0 - 0 = 0 \\
            W2 &= 3 - 1 = 2 \\
            W3 &= 8 - 5 = 3 \\
            W4 &= 16 - 16 = 0 \\
        \end{align*}
        
        \item \textbf{Step 3: Calculate the average:}
        \[
            \text{Average Waiting Time} = \frac{0 + 2 + 3 + 0}{4} = 1.25
        \]
    \end{itemize}
\end{itemize}

\subsection{Variation 5.3}
\begin{itemize}
    \item \textbf{Question:} Consider the following processes with their burst times and arrival times:
    
    \begin{tabular}{|l|l|l|}
        \hline
        \textbf{Process} & \textbf{Arrival Time (ms)} & \textbf{Burst Time (ms)} \\
        \hline
        P1 & 0 & 8 \\
        P2 & 2 & 5 \\
        P3 & 5 & 6 \\
        P4 & 20 & 9 \\
        \hline
    \end{tabular}
    
    Using the \textbf{First-Come-First-Served (FCFS)} scheduling algorithm, what is the \textbf{average waiting time}?
    
    \begin{tabular}{|l|l|}
        \hline
        \textbf{Correct Answer(s)} & \textbf{Incorrect Answers} \\
        \hline
        3.5 & 7 \\
         & 3.75 \\
         & 3.25 \\
         & 6.25 \\
        \hline
    \end{tabular}
    
    \item \textbf{Solution \& Feedback:}
    \begin{itemize}
        \item \textbf{Step 1: Calculate the finish time for each process:}
        \begin{align*}
            \text{Finish Time (P1)} &= \text{Arrival Time (P1)} + \text{Burst Time (P1)} = 0 + 8 = 8 \\
            \text{Finish Time (P2)} &= \max(\text{Finish Time (P1)}, \text{Arrival Time (P2)}) + \text{Burst Time (P2)} = 8 + 5 = 13 \\
            \text{Finish Time (P3)} &= \max(\text{Finish Time (P2)}, \text{Arrival Time (P3)}) + \text{Burst Time (P3)} = 13 + 6 = 19 \\
            \text{Finish Time (P4)} &= \max(\text{Finish Time (P3)}, \text{Arrival Time (P4)}) + \text{Burst Time (P4)} = 20 + 9 = 29 \\
        \end{align*}
        
        \item \textbf{Step 2: Calculate the waiting time (W) for each process: W = Start Time – Arrival Time}
        \begin{align*}
            W1 &= 0 - 0 = 0 \\
            W2 &= 8 - 2 = 6 \\
            W3 &= 13 - 5 = 8 \\
            W4 &= 20 - 20 = 0 \\
        \end{align*}
        
        \item \textbf{Step 3: Calculate the average:}
        \[
            \text{Average Waiting Time} = \frac{0 + 6 + 8 + 0}{4} = 3.5
        \]
    \end{itemize}
\end{itemize}

\subsection{Variation 6.1}
\begin{itemize}
    \item \textbf{Question:} What is the complexity of the following algorithm?
    
    \begin{lstlisting}
    def compute_result(arr):
        n = len(arr)
        result = 0
        for i in range(n):
            for j in range(i + 1, n):
                result += arr[i] * arr[j]
        for k in range(n):
            result += arr[k]
        return result
    \end{lstlisting}
    
    \begin{tabular}{|l|l|}
        \hline
        \textbf{Correct Answer(s)} & \textbf{Incorrect Answers} \\
        \hline
        O(n\textsuperscript{2}) & O(log(n)) \\
         & O(n) \\
         & O(n log(n)) \\
         & O(n\textsuperscript{3}) \\
        \hline
    \end{tabular}
    
    \item \textbf{Solution \& Feedback:}
    \begin{itemize}
        \item The total number of instructions is \( 1 + 1 + n * n/2 + n + 1 = (1/2)n^2 + n + 3 \).
        \item The quadratic term \((1/2)n^2\) dominates, as it grows faster than the linear term (n) or the constant term (3).
        \item Therefore, ignoring the constant factor \((1/2)\), the algorithm is in the O(n\textsuperscript{2}) complexity class.
    \end{itemize}
\end{itemize}

\subsection{Variation 6.2}
\begin{itemize}
    \item \textbf{Question:} What is the complexity of the following algorithm?
    
    \begin{lstlisting}
    def compute_result(arr):
        n = len(arr)
        result = 0
        for i in range(n):
            for j in range(i + 1, n):
                result += arr[i] * arr[j]
        for k in range(n):
            result += arr[k]
        return result
    \end{verbatim}
    
    \begin{tabular}{|l|l|}
        \hline
        \textbf{Correct Answer(s)} & \textbf{Incorrect Answers} \\
        \hline
        O(n\textsuperscript{2}) & O(log(n)) \\
         & O(n) \\
         & O(n log(n)) \\
         & O(n\textsuperscript{3}) \\
        \hline
    \end{lstlisting}
    
    \item \textbf{Solution \& Feedback:}
    \begin{itemize}
        \item The total number of instructions is \( 1 + 1 + n * (n/2 + n) + 1 = (3/2)n^2 + 3 \).
        \item The quadratic term \((3/2)n^2\) dominates, as it grows faster than the constant term (3).
        \item Therefore, ignoring the constant factor \((3/2)\), the algorithm is in the O(n\textsuperscript{2}) complexity class.
    \end{itemize}
\end{itemize}

\subsection{Variation 6.3}
\begin{itemize}
    \item \textbf{Question:} What is the complexity of the following algorithm?
    
    \begin{lstlisting}
    def compute_result(arr):
        n = len(arr)
        result = 0
        for i in range(n):
            for j in range(i + 1, n):
                result += arr[i] * arr[j]
                for k in range(n):
                    result += arr[k]
        return result
    \end{lstlisting}
    
    \begin{tabular}{|l|l|}
        \hline
        \textbf{Correct Answer(s)} & \textbf{Incorrect Answers} \\
        \hline
        O(n\textsuperscript{3}) & O(log(n)) \\
         & O(n) \\
         & O(n log(n)) \\
         & O(n\textsuperscript{2}) \\
        \hline
    \end{tabular}
    
    \item \textbf{Solution \& Feedback:}
    \begin{itemize}
        \item The total number of instructions is \( 1 + 1 + n * (n/2 + n^2) = n^3 + (1/2)n^2 + 2 \).
        \item The cubic term \( n^3 \) dominates, as it grows faster than the quadratic term \((1/2)n^2\) and constant term (2).
        \item Therefore, the algorithm is in the O(n\textsuperscript{3}) complexity class.
    \end{itemize}
\end{itemize}

\subsection{Variation 7.1}
\begin{itemize}
    \item \textbf{Question:} You have written a program that has a fixed startup time, and it uses an algorithm which you know is O(n\textsuperscript{3}). You have tested this program with 1 data item, and it takes 150 ms to execute. With 1000 data items it takes 450 ms to run. Estimate how long (in ms) this program will take for 4000 data items?
    
    \begin{tabular}{|l|l|}
        \hline
        \textbf{Correct Answer(s)} & \textbf{Incorrect Answers} \\
        \hline
        19350 ms & 24450 ms (taking 3\textsuperscript{4} instead of 4\textsuperscript{3}) \\
         & 28800 ms (no fixed startup cost i.e. 64 x 450) \\
         & 1350 ms (fixed + 4 x 300) \\
         & 1800 ms (4 times the cost of 1000 i.e. 4 x 450) \\
        \hline
    \end{tabular}
    
    \item \textbf{Solution \& Feedback:}
    \begin{itemize}
        \item With 1 item it takes 150 ms. This is approx. the fixed cost of the algorithm.
        \item With n=1000 it takes 450 ms: 150 ms startup time + 300 ms to process the 1000 items.
        \item The main algorithm is O(n\textsuperscript{3}). Therefore, if we have 4 times as much data i.e. 4000 items, we can estimate that it will take 4\textsuperscript{3} = 64 times as long to execute that part of the program + the fixed startup time.
        \item Therefore, it will take \( 150 + 64 \times 300 = 150 + 19200 = 19350 \) ms.
    \end{itemize}
\end{itemize}

\subsection{Variation 7.2}
\begin{itemize}
    \item \textbf{Question:} You have written a program that has a fixed startup time, and it uses an algorithm which you know is O(n\textsuperscript{4}). You have tested this program with 1 data item, and it takes 120 ms to execute. With 500 data items it takes 300 ms to run. Estimate how long (in ms) this program will take for 1500 data items?
    
    \begin{tabular}{|l|l|}
        \hline
        \textbf{Correct Answer(s)} & \textbf{Incorrect Answers} \\
        \hline
        14700 ms & 11640 ms (taking 4\textsuperscript{3} instead of 3\textsuperscript{4}) \\
         & 24300 ms (no fixed startup cost i.e. 81 x 300) \\
         & 660 ms (fixed + 3 x 180) \\
         & 900 ms (3 times the cost of 500 i.e. 3 x 300) \\
        \hline
    \end{tabular}
    
    \item \textbf{Solution \& Feedback:}
    \begin{itemize}
        \item With 1 item it takes 120 ms. This is approx. the fixed cost of the algorithm.
        \item With n=500 it takes 300 ms: 120 ms startup time + 180 ms to process the 500 items.
        \item The main algorithm is O(n\textsuperscript{4}). Therefore, if we have 3 times as much data i.e. 1500 items, we can estimate that it will take 3\textsuperscript{4} = 81 times as long to execute that part of the program + the fixed startup time.
        \item Therefore, it will take \( 120 + 81 \times 180 = 120 + 14580 = 14700 \) ms.
    \end{itemize}
\end{itemize}

\subsection{Variation 7.3}
\begin{itemize}
    \item \textbf{Question:} You have written a program that has a fixed startup time, and it uses an algorithm which you know is O(n\textsuperscript{5}). You have tested this program with 1 data item, and it takes 80 ms to execute. With 700 data items it takes 230 ms to run. Estimate how long (in ms) this program will take for 2100 data items?
    
    \begin{tabular}{|l|l|}
        \hline
        \textbf{Correct Answer(s)} & \textbf{Incorrect Answers} \\
        \hline
        36530 ms & 18830 ms (taking 5\textsuperscript{3} instead of 3\textsuperscript{5}) \\
         & 55890 ms (no fixed startup cost i.e. 243 x 230) \\
         & 530 ms (fixed + 3 x 150) \\
         & 690 ms (3 times the cost of 230 i.e. 3 x 230) \\
        \hline
    \end{tabular}
    
    \item \textbf{Solution \& Feedback:}
    \begin{itemize}
        \item With 1 item it takes 80 ms. This is approx. the fixed cost of the algorithm.
        \item With n=700 it takes 230 ms: 80 ms startup time + 150 ms to process the 700 items.
        \item The main algorithm is O(n\textsuperscript{5}). Therefore, if we have 3 times as much data i.e. 2100 items, we can estimate that it will take 3\textsuperscript{5} = 243 times as long to execute that part of the program + the fixed startup time.
        \item Therefore, it will take \( 80 + 243 \times 150 = 80 + 36450 = 36530 \) ms.
    \end{itemize}
\end{itemize}
\subsection{Variation 8.1}
\textbf{Question}\\
Consider a processor with 3 levels of memory, Level 1 Cache, Level 2 Main Memory (RAM) and Level 3 Secondary Storage (Disk) as represented in the figure below:\\
\begin{itemize}
    \item \textbf{Level 1 (Cache):} Access time = 5 ns, Hit ratio = 0.85
    \item \textbf{Level 2 (RAM):} Access time = 100 ns, Hit ratio = 0.75
    \item \textbf{Level 3 (Disk):} Access time = 4 ms
\end{itemize}
Level 1 (Cache) has an access time of 5ns, Level 2 (RAM) has an access time of 100ns, and Level 3 (Disk) has an access time of 4ms. The hit-ratios for Cache and RAM are 0.85 \& 0.75, respectively. \\

If a data item is not found at one level, it is fetched from the next level and the reference restarts. 

\noindent \textbf{What is the average time (in nanoseconds) required to access a referenced word in this system?}

\textbf{Options}
\begin{tabular}{ll}
    \textbf{Correct Answer(s)} & \textbf{Incorrect Answers} \\
    \hline
    150,020 ns & 170 ns (assuming 4 ms = 4000 ns) \\
    & 1,000,109.25 ns (using incorrect probabilities) \\
    & 1,000,079.25 ns (using incorrect probabilities \& access times) \\
    & 150,015.50 ns (using incorrect access times) \\
\end{tabular}

\textbf{Solution \& Feedback}
The average memory access time (AMAT) is calculated as:
\[
\text{AMAT} = (\text{Hit Ratio}_1 \times \text{Time}_1) + (\text{Miss Ratio}_1 \times \text{Hit Ratio}_2 \times \text{Time}_2) + (\text{Miss Ratio}_1 \times \text{Miss Ratio}_2 \times \text{Time}_3)
\]
where:
\begin{itemize}
    \item $\text{Time}_1$ (Cache hit): 5 ns
    \item $\text{Time}_2$ (RAM access after cache miss): 5 ns + 100 ns = 105 ns
    \item $\text{Time}_3$ (Disk access after RAM miss): 5 ns + 100 ns + 4 ms = 4,000,105 ns
\end{itemize}

\noindent Substituting the values:
\[
0.85 \times 5 + 0.15 \times 0.75 \times 105 + 0.15 \times 0.25 \times 4,000,105 = 150,020 \text{ ns}
\]

\noindent \textbf{Key Points:}
\begin{itemize}
    \item The disk access time (4 ms) dominates the average time due to its magnitude.
    \item The hit/miss ratios must be applied sequentially.
    \item All access times include the overhead of restarting the reference from the top level.
\end{itemize}



% ===== Question 8.2 =====
\subsection{Variation 8.2}
\textbf{Question}
Consider a processor with 3 levels of memory:
\begin{itemize}[leftmargin=*]
    \item \textbf{Level 1 (Cache):} Access time = 7 ns, Hit ratio = 0.75
    \item \textbf{Level 2 (RAM):} Access time = 150 ns, Hit ratio = 0.90
    \item \textbf{Level 3 (Disk):} Access time = 5 ms
\end{itemize}

\noindent Calculate the average access time (in ns) for a referenced word.

\textbf{Options}
\begin{tabular}{ll}
    \textbf{Correct Answer} & \textbf{Incorrect Answers} \\
    \hline
    125,044.50 ns & 169.50 ns (assuming 5 ms = 5000 ns) \\
    & 500,162.25 ns (incorrect probabilities) \\
    & 500,140.25 ns (incorrect probabilities/access times) \\
    & 125,039 ns (incorrect access times) \\
\end{tabular}

\textbf{Solution}
\[
\begin{aligned}
\text{AMAT} &= 0.75 \times 7 + 0.25 \times 0.90 \times (7 + 150) \\
           &+ 0.25 \times 0.10 \times (7 + 150 + 5,000,000) \\
           &= 125,044.50 \text{ ns}
\end{aligned}
\]

% ===== Question 8.3 =====
\subsection{Variation 8.3}
\textbf{Question}
Memory hierarchy with:
\begin{itemize}[leftmargin=*]
    \item \textbf{Cache:} 10 ns, Hit ratio = 0.80
    \item \textbf{RAM:} 200 ns, Hit ratio = 0.60
    \item \textbf{Disk:} 3 ms
\end{itemize}

\textbf{Solution}
\[
0.80 \times 10 + 0.20 \times 0.60 \times 210 + 0.20 \times 0.40 \times 3,000,210 = 240,050 \text{ ns}
\]

% ===== Question 9.1 =====
\subsection{Variation 9.1}
\textbf{Question}
What is the complexity of the following algorithm:

\begin{lstlisting}[language=Python]
def some_fun(n):
    result = 0
    if n < 0:
        for i in range(abs(n)):
            for j in range(i + 1, n):
                result += i*j
                for k in range(n):
                    result += k
    else:
        result = recursive_fun(n)
    return result

def recursive_fun(n):
    if n <= 1:
        return 1
    else:
        count = 0
        for i in range(n):
            count += recursive_fun(i)
        return count
\end{lstlisting}

\begin{center}
\begin{tabular}{|l|l|}
\hline
\textbf{Correct Answer(s)} & \textbf{Incorrect Answers} \\
\hline
$O(2^n)$ & $O(n^2) O(nlogn) O(n^3)$ \\
\hline
\end{tabular}
\end{center}

\textbf{Solution \& Feedback}
In the some\_fun function, the time complexity for $n<0$ is $O(n^3)$, but the recursive\_fun function has a much worse time complexity as explained below:

\textbf{Base Case:} When n<=1, the function returns a constant value (O(1)).
\begin{itemize}
\item T(0) = O(1)
\item T(1) = O(1)
\end{itemize}

\textbf{Recursive Case:} For $n>1$, each iteration of the loop, the function calls itself with input i:
\begin{itemize}
\item At T(2), we make $2^2-1=3$ recursive calls.
\item At T(3), we make $2^3-1=7$ recursive calls.
\end{itemize}

So, in general, $T(n)=2^n-1$. Hence the overall time complexity is $O(2^n)$.

% ===== Question 9.2 =====
\subsection{Variation 9.2}
\textbf{Question}
What is the complexity of the following algorithm:

\begin{lstlisting}[language=Python]
def some_fun(n):
    result = 0
    if n < 0:
        for i in range(abs(n)):
            for j in range(n):
                result += i*j
    else:
        result = recursive_fun(n)
    return result

def recursive_fun(n):
    if n <= 1:
        return 1
    else:
        count = 0
        for i in range(n):
            count += recursive_fun(i)
        return count
\end{lstlisting}

\begin{center}
\begin{tabular}{|l|l|}
\hline
\textbf{Correct Answer(s)} & \textbf{Incorrect Answers} \\
\hline
$O(2^n)$ & $O(n^2)$ $O(nlogn)$ $O(n^3)$ \\
\hline
\end{tabular}
\end{center}

\textbf{Solution \& Feedback}
In the some\_fun function, the time complexity for n<0 is $O(n^2)$, but the recursive\_fun function has a much worse time complexity as explained below:

\textbf{Base Case:} When n<=1, the function returns a constant value (O(1)).
\begin{itemize}
\item T(0) = O(1)
\item T(1) = O(1)
\end{itemize}

\textbf{Recursive Case:} For n>1, each iteration of the loop, the function calls itself with input i:
\begin{itemize}
\item At $T(2)$, we make $2^2-1=3$ recursive calls.
\item At $T(3)$, we make $2^3-1=7$ recursive calls.
\end{itemize}

So, in general, $T(n)=2^n-1$. Hence the overall time complexity is $O(2^n)$.

% ===== Question 9.3 =====
\subsection{Variation 9.3}
\textbf{Question}
What is the complexity of the following algorithm:

\begin{lstlisting}[language=Python]
def some_fun(n):
    result = 0
    if n < 0:
        for i in range(abs(n)):
            result += i
    else:
        result = recursive_fun(n)
    return result

def recursive_fun(n):
    if n <= 1:
        return 1
    else:
        count = 0
        for i in range(n):
            count += recursive_fun(i)
        return count
\end{lstlisting}

\begin{center}
\begin{tabular}{|l|l|}
\hline
\textbf{Correct Answer(s)} & \textbf{Incorrect Answers} \\
\hline
$O(2^n)$ & $O(n^2)$ $O(nlogn)$ $O(n^3)$ \\
\hline
\end{tabular}
\end{center}

\textbf{Solution \& Feedback}
In the some\_fun function, the time complexity for $n<0$ is O(n), but the recursive\_fun function has a much worse time complexity as explained below:

\textbf{Base Case:} When $n<=1$, the function returns a constant value (O(1)).
\begin{itemize}
\item T(0) = O(1)
\item T(1) = O(1)
\end{itemize}

\textbf{Recursive Case:} For n>1, each iteration of the loop, the function calls itself with input i:
\begin{itemize}
\item At T(2), we make $2^2$-1=3 recursive calls.
\item At T(3), we make $2^3$-1=7 recursive calls.
\end{itemize}

So, in general, \(T(n)=2^n-1\). Hence the overall time complexity is $O(2^n)$.

% ===== Question 10.1 =====
\subsection{Variation 10.1}
\textbf{Question}\\
Consider the following table that shows the history of CPU burst times for four processes (P1, P2, P3, and P4). The CPU burst times are measured in milliseconds, and exponential averaging is used to predict the next CPU burst. The initial prediction ($T_0$) for each process is 10 ms, and the value of $\alpha$ is 0.25.\\
\begin{center}
\begin{tabular}{|c|c|c|c|c|}
\hline
Process & Actual Burst ($t_0$) & Predicted Burst ($T_1$) & Actual Burst ($t_1$) & Predicted Burst ($T_2$) \\
\hline
P1 & 14 & ? & 7 & ? \\
P2 & 6 & ? & 9 & ? \\
P3 & 2 & ? & 4 & ? \\
P4 & 18 & ? & 16 & ? \\
\hline
\end{tabular}
\end{center}

Calculate the next predicted CPU burst times ($T_1$ and $T_2$) for each process using exponential averaging and select the correct options from below.

\textbf{Note:} The options are listed in Px [$T_1$, $T_2$] format.

\begin{center}
\begin{tabular}{|l|l|}
\hline
\textbf{Correct Answer(s)} & \textbf{Incorrect Answers} \\
\hline
P1 [11, 10]; P2 [9, 9]; P3 [8, 7]; P4 [12, 13] & P1 [11, 9.25]; P2 [9, 9.75]; P3 [8, 8.5]; P4 [12, 11.5] \\
 & P1 [13, 8.5]; P2 [7, 8.5]; P3 [4, 4]; P4 [16, 16] \\
 & P1 [13, 7.75]; P2 [7, 9.25]; P3 [4, 5.5]; P4 [16, 14.5] \\
\hline
\end{tabular}
\end{center}

\textbf{Solution \& Feedback}
We compute the predicted CPU burst times using:
\[ T_{n+1} = \alpha \times t_n + (1 - \alpha) \times T_n \]

\begin{center}
\begin{tabular}{|c|c|c|c|c|}
\hline
Process & Actual Burst ($t_0$) & Predicted Burst ($T_1$) & Actual Burst ($t_1$) & Predicted Burst ($T_2$) \\
\hline
P1 & 14 & \(14 \times 0.25 + 10 \times 0.75 = 11\) & 7 & \(7 \times 0.25 + 11 \times 0.75 = 10 \) \\
P2 & 6 & 9 & 9 & 9 \\
P3 & 2 & 8 & 4 & 7 \\
P4 & 18 & 12 & 16 & 13 \\
\hline
\end{tabular}
\end{center}

% ===== Question 10.2 =====
\subsection{Variation 10.2}
\textbf{Question}
Consider the following table with initial prediction ($T_0$) = 5 ms and $\alpha$ = 0.40:

\begin{center}
\begin{tabular}{|c|c|c|c|c|}
\hline
Process & Actual Burst ($t_0$) & Predicted Burst ($T_1$) & Actual Burst ($t_1$) & Predicted Burst ($T_2$) \\
\hline
P1 & 10 & ? & 12 & ? \\
P2 & 15 & ? & 24 & ? \\
P3 & 5 & ? & 10 & ? \\
P4 & 20 & ? & 16 & ? \\
\hline
\end{tabular}
\end{center}

\begin{center}
\begin{tabular}{|l|l|}
\hline
\textbf{Correct Answer(s)} & \textbf{Incorrect Answers} \\
\hline
P1 [7, 9]; P2 [9, 15]; P3 [5, 7]; P4 [11, 13] & P1 [7, 7.8]; P2 [9, 12.6]; P3 [5, 7]; P4 [11, 9.4] \\
 & P1 [8, 10.4]; P2 [11, 18.8]; P3 [5, 8]; P4 [14, 15.2] \\
 & P1 [8, 9.2]; P2 [11, 16.4]; P3 [5, 8]; P4 [14, 11.6] \\
\hline
\end{tabular}
\end{center}

\textbf{Solution \& Feedback}
Calculations with $\alpha$ = 0.40:

\begin{center}
\begin{tabular}{|c|c|c|c|c|}
\hline
Process & $t_0$ & $T_1$ & $t_1$ & $T_2$ \\
\hline
P1 & 10 & 7 & 12 & 9 \\
P2 & 15 & 9 & 24 & 15 \\
P3 & 5 & 5 & 10 & 7 \\
P4 & 20 & 11 & 16 & 13 \\
\hline
\end{tabular}
\end{center}

% ===== Question 10.3 =====
\subsection{Variation 10.3}
\textbf{Question}
Consider the following table with initial prediction ($T_0$) = 20 ms and $\alpha$ = 0.60:

\begin{tabular}{|c|c|c|c|c|}
\hline
Process & Actual Burst ($t_0$) & Predicted Burst ($T_1$) & Actual Burst ($t_1$) & Predicted Burst ($T_2$) \\
\hline
P1 & 8 & ? & 14 & ? \\
P2 & 10 & ? & 16 & ? \\
P3 & 15 & ? & 12 & ? \\
P4 & 5 & ? & 6 & ? \\
\hline
\end{tabular}

\begin{center}
\begin{tabular}{|l|l|}
\hline
\textbf{Correct Answer(s)} & \textbf{Incorrect Answers} \\
\hline
P1 [12.8, 13.52]; P2 [14, 15.2]; P3 [17, 14]; P4 [11, 8] & P1 [12.8, 16.4]; P2 [14, 17.6]; P3 [17, 15.2]; P4 [11, 11.6] \\
 & P1 [15.2, 14.72]; P2 [16, 16]; P3 [18, 15.6]; P4 [14, 10.8] \\
 & P1 [15.2, 17.6]; P2 [16, 18.4]; P3 [18, 16.8]; P4 [14, 14.4] \\
\hline
\end{tabular}
\end{center}

\textbf{Solution \& Feedback}\\
Calculations with $\alpha$ = 0.60:
\begin{tabular}{|c|c|c|c|c|}
\hline
Process & $t_0$ & $T_1$ & $t_1$ & $T_2$ \\
\hline
P1 & 8 & 12.8 & 14 & 13.52 \\
P2 & 10 & 14 & 16 & 15.2 \\
P3 & 15 & 17 & 12 & 14 \\
P4 & 5 & 11 & 6 & 8 \\
\hline
\end{tabular}